% !TEX TS-program = lualatex+se
% !TEX encoding = UTF-8
\documentclass[12pt, a4paper]{book}

% --- USTAWIENIA STRONY ---
\usepackage{geometry}
\geometry{
%	showframe=true,
%	showcrop=true,
	left=20mm,
	right=20mm,
	top=20mm,
	bottom=18mm,
	marginparwidth=0mm,
	marginparsep=0mm
	}
%\usepackage[letterpaper,showframe]{geometry}	
% usuwanie nagłówków/stopek z\,dokumentu
\usepackage[indent=0pt]{parskip} % ustawianie paragrafów
%\setstretch{1.5} % półtorej linii odległości między wierszami
\usepackage{titlesec} % ustawianie nagłówków
%\usepackage{fullpage} % redukcja marginesów
\usepackage{setspace} % odstępy między liniami
\usepackage{enumitem}% listy numerowane
% -----------------------------

%----------- style nagłówków ----------------
\newcommand{\headingone}{\scshape\Large}
\newcommand{\headingtwo}{\large}
\newcommand{\margintext}[1]{\textit{\small #1}}
% przesunięcia poziome
\newcommand{\smallwidth}{\hskip 1em}
\newcommand{\medwidth}{\hskip 3pt}
\newcommand{\largewidth}{\hskip 12pt}
% przesunięcia pionowe
\newcommand{\linespace}{\vskip 6pt}
\newcommand{\paragraphspace}{\vskip 10pt}
\newcommand{\segmentspace}{\vskip 16pt}
\newcommand{\sectionspace}{\vskip 28pt}
%------------------------------------------

% --- GREGORIO ---	
\usepackage{gregoriotex}
\gresetlinecolor{gregoriocolor}
\gresetheadercapture{annotation}{greannotation}{string}% odczytuje adnotacje nad inicjałem jako ciąg, co pozwala na jego renderowanie
\gresetheadercapture{commentary}{grecommentary}{string}% odczytuje komentarz jako ciąg, co pozwala na jego renderowanie
\grechangedim{commentaryraise}{0.3cm}{scalable} % podnosi komentarz wystarczająco ponad górną linię
\grechangestyle{annotation}{\color{black}\small}% format adnotacji
\grechangestyle{commentary}{\color{gregoriocolor}\small}% format komentarza
%\grechangestyle{commentary}{\color{gregoriocolor}\small\itshape}% format komentarza
\grechangestyle{abovelinestext}{\color{gregoriocolor}\small\itshape}% format napisu nad liniami
\grechangestyle{initial}{\fontsize{36}{36}\selectfont} % zmniejsza wielkość inicjału do bardziej proporcjonalnego rozmiaru (=36pt)
\grechangedim{intersyllablespacenotes}{0.21cm}{scalable}% Minimalna odległość między nutami różnych sylab.

% umieszczenie zawartości składu ich unikalnego argumentu po lewej i prawej stronie każdej linii
\let\grelocalleftbox\localleftbox
\let\grelocalrightbox\localrightbox

\newcommand{\wstawplik}[1]{\input{/Users/Laser/Documents/Latex/include/#1}}%skraca kod
%\newcommand{\separator}{\greseparator{3}{40}}
%\newfontface\vectis{VectisMonumentalW01-Regular}
%\def\greinitialformat#1{\fontsize{20}{20}\vectis #1}
% ------------------

%\usepackage[hyperindex=true]{hyperref}% przypisy dolne

%---- usun konflikt między paracol i luacolor ----
\wstawplik{conparlua}
%-----------------------------------------------

% --- SPIS TREŚCI -------
%\usepackage{imakeidx}
%\makeindex[title=Alfabetyczny spis program=texindy,options=-M styl.xdy]
%\usepackage{filecontents}
%\begin{filecontents*}{styl.xdy}
%;;; xindy style file
%(markup-locclass-list :open "\dotfill" :sep "")
%(define-attributes (( "textbf" "default" )) )
%(markup-locref   :attr  "textbf"     :open "\textbf{\hyperpage{" :close "}}")
%(markup-locref   :attr  "textit"     :open "\textit{\hyperpage{" :close "}}")
%(markup-locref   :attr  "textttt"     :open "\textttt{\hyperpage{" :close "}}")
%(markup-locref   :attr  "texttsc"     :open "\texttsc{\hyperpage{" :close "}}")
%(markup-locref   :attr  "default"     :open "\hyperpage{" :close "}")
%\end{filecontents*}
% -------------------------

% ---  Konfiguracja kolumn ---
\usepackage{paracol}
\setcolumnwidth{.32\textwidth, .32\textwidth, .32\textwidth}
\setlength{\columnseprule}{0.4pt}
\setlength{\columnsep}{.02\textwidth}
\colseprulecolor{red}% wymagany pakiet xcolor
\footnotelayout{m}
% ----------------------------

\makeatletter 
\patchcmd\pcol@hfil
 {\copy\pcol@tempboxa}
 {\luacolorresetboxcolor\pcol@tempboxa\copy\pcol@tempboxa}{}{\fail}
\makeatother
% ----------------------------

%% --- UKŁAD DWUKOLUMNOWY Z TŁUMACZENIEM ---
%\usepackage{tcolorbox}
%\newtcolorbox{nutybox}{size=minimal, colback=white, colframe=red, arc=0mm, width=\linewidth, halign=left, valign=center}
%\newtcolorbox{tlumbox}{size=minimal, colback=white, colframe=red, arc=0mm, width=\linewidth, halign=left, valign=top, fontupper=\small\itshape} 
%\newcommand{\nuty}[1]{\begin{nutybox}#1 \end{nutybox}}
%\newcommand{\tlum}[1]{\begin{tlumbox}\tlumacz{#1} \end{tlumbox}}
%% ------------------------------------------------------

\usepackage{Tabbing}
\usepackage{longtable}
\usepackage[savepos]{zref}

% --- CZCIONKI DOKUMENTU ---
\usepackage[latin,polish]{babel}
\usepackage{fontspec}
%\usepackage{times}
%\usepackage[math]{anttor}
%\usepackage{tgtermes}
%\usepackage{mathptmx}
\usepackage{ebgaramond}
\usepackage{libertine}
\defaultfontfeatures{Ligatures=TeX}

%\setmainfont{TeX Gyre Termes}
%\setmainfont{TeX Gyre Heros}
%\setmainfont{Calluna Regular}
\setmainfont{EB Garamond}
%\setmainfont{Bodoni Moda}
%\newfontfamily{\sekfont}{Times New Roman}
\newfontfamily{\sekfont}{Libertinus Serif}
\newcommand\sekcjafont{\fontsize{16}{16}\selectfont \sekfont}

%----- Strona tytułowa -------------
\newfontfamily{\tytfont}{JSL Ancient}
\newcommand\tytulfont[1]{\fontsize{36}{30}\selectfont \tytfont#1\par}
\newcommand\podtytulfont[1]{\fontsize{18}{10}\selectfont \tytfont#1\par}

%\setmainfont[<font features>]{<name of font>} % sets the roman font
%\setsansfont[<font features>]{<name of font>} % sets the sans font
%\setmonofont[<font features>]{<name of font>} % sets the monospace font

% --- KONIEC CZCIONKI DOKUMENTU ---

% --- GRAFIKA I KOLORY ---
\usepackage{xcolor} % kolory
\usepackage{graphicx} % umożliwia import obrazów
\graphicspath{{img/}} % katalog obrazów
\usepackage{pgfornament}% ornamenty graficzne
% --- KONIEC GRAFIKA I KOLORY -----------

% --- KOMENDY WŁASNE ---
\newcommand{\red}[1]{\textcolor{red}{#1}}
\newcommand{\black}[1]{\textcolor{black}{#1}}
% --------------------------------
\newcommand{\tlumacz}[1]{\textit{#1}}
\newcommand{\inicjal}[2]{\textcolor{red}{\LARGE\textbf{#1}}{#2}}
%--------------------------------
\titlespacing*{\chapter}{0pt}{0pt}{5pt}	
\titleformat{\chapter}[hang]{\filcenter\LARGE\scshape\bfseries}{}{0pt}{} % format nagłówków
\titlespacing*{\section}{0pt}{0pt}{0pt}	
\titleformat{\section}[hang]{\filcenter\large\scshape}{}{0pt}{} % format nagłówków

% ----------- RUBRYKA -------------------
\newcommand{\rubryka}[1]{
	\begin{spacing}{0.8}
	\begin{center}
		\parbox{0.9\textwidth}{\centering\textcolor{red}{\small\textit{#1}}}
	\end{center}
	\end{spacing}
	\linespace}
% ----------- SEKCJA --------------------
% -----------------------------------------
%\renewcommand\fbox{\fcolorbox{red}{white}}% czerwona ramka fbox-a
\newcommand{\sekcja}[2]{
	\begin{spacing}{0.6}
	\begin{center}
		\section*{#1}
		\linespace
		\begin{minipage}{\linewidth}
			\centering\red{\small\textit{#2}}
		\end{minipage}
	\end{center}	
	\end{spacing}
	\paragraphspace}
%------------DEKOR ---------------------
\newcommand{\dekor}[2]{
	\begin{center}
		\vspace*{\fill}
		\includegraphics[width=#1\textwidth]{#2}
		\vspace*{\fill}
	\end{center}}
%------- Zarządzanie wersami w psalmach z tlumaczeniem -------------------------
\newcounter{licznik}
\newcommand{\licznikZero}{\setcounter{licznik}{0}}% zeruj licznik
\newcommand{\licznikUp}{\refstepcounter{licznik}}% zwiększ licznik 
\licznikZero
\newcommand{\BeginParallel}{\begin{spacing}{0.9} \begin{paracol}{2} \licznikZero}
\newcommand{\lat}[1]{\selectlanguage{latin}\ensurevspace{2\baselineskip}\begin{minipage}[t]{\linewidth}\ifnum \value{licznik}>0 \licznikUp \thelicznik .~{#1}\par \else \licznikUp{#1}\par \fi \end{minipage}\switchcolumn}
\newcommand{\latN}[1]{\selectlanguage{latin}\ensurevspace{2\baselineskip}\begin{minipage}[t]{\linewidth}{#1}\par\end{minipage}\switchcolumn}% kolumna bez numeracji wersów
\newcommand{\pol}[1]{\selectlanguage{polish}\begin{minipage}[t]{\linewidth}{#1}\par\end{minipage}\switchcolumn*}
\newcommand{\EndParallel}{\end{paracol} \end{spacing}}

% --- ZNAKI SPECJALNE ---
% 
%\newcommand <c>†</c>
% --------------------------

\begin{document}
\makeatletter
\grechangedim{interwordspacetext}{\unexpanded{\the\fontdimen2\font plus \the\gluestretch\gre@space@skip@interwordspacenotes minus \the\glueshrink\gre@space@skip@interwordspacenotes}}{fixed}%
\makeatother

\pagestyle{plain}	
\large
\begin{spacing}{0.7}
\begin{center}
	\section*{Te Deum laudamus (tonus solemnis)}
	\small\textit{Liber Usualis, 1961, p. 1832}
\end{center}	
\end{spacing}
\vspace{7mm}
\setlength{\parindent}{0pt}
\gresetbracerendering{font}

\pagenumbering{gobble} % nie jest wyprowadzany żaden numer strony, chociaż numer jest nadal resetowany. Odniesienia do tej strony również są puste.
\greannotation{Hymn.} \greannotation{III.} \gabcsnippet{(c4) TE(e) De(gh)um(h) lau(hg)dá(h!iwj)mus :(ih..) *(;) te(e) Dó(gh)mi(h)num(h') con(g)fi(hi)té(ivHG)mur.(g.) (::) 
Te(gi) ae(h)tér(j')num(i) Pa(hi)trem(h.) (;) o(e)mnis(gh) ter(h)ra(h') ve(g)ne(hi)rá(ivHG)tur.(g.) (::) 
Ti(gi)bi(h) o(j')mnes(i) An(hi)ge(hg)li,(h.) (;) ti(e)bi(gh) Cae(h)li(h) et(h) u(h)ni(h)vér(h)sae(h') pot(g)e(hi)stá(ivHG)tes :(g.) (::) 
Ti(gi)bi(h) Ché(j')ru(i)bim(h) et(h) Sé(hi)ra(hg)phim(h.) (;) in(e)ces(gh)sá(h)bi(h)li(h) vo(h')ce(g) pro(hi)clá(ivHG)mant :(g.) (::) 
<alt>Fit reverentia</alt> 
San(g.h!iwj)ctus :(ih..) (::) 
San(g.h!iwj)ctus :(ih..) (::) 
San(eg!h!iwj)ctus(ih..) (;) Dó(h)mi(h)nus(h') De(g)us(hi) Sá(i)ba(hg)oth.(g.) (::) 
Ple(gi)ni(i) sunt(h) cae(j')li(i) et(h) ter(hi)ra(h.) (;) ma(e)je(gh)stá(h)tis(h) gló(h')ri(g)ae(hi) tu(ivHG)ae.(g.) (::) 
Te(gih) glo(j')ri(i)ó(hi)sus(h.) (;) A(h)po(h)sto(h')ló(g)rum(hi) cho(ivHG)rus :(g.) (::) 
Te(gih) Pro(j')phe(i)tá(hi)rum(h.) (;) lau(h)dá(h')bi(g)lis(hi) nú(i)me(hg)rus :(g.) (::) 
Te(gih) Már(j')ty(i)rum(h') can(h)di(h)dá(hi)tus(h.) (;) lau(h')dat(g) ex(hi)ér(i)ci(hg)tus.(g.) (::) 
Te(gi) per(h) or(j')bem(i) ter(h)rá(hi)rum(h.) (;) san(e)cta(gh) con(h)fi(h)té(h')tur(g) Ec(hi)clé(i)si(hg)a :(g.) (::) 
Pa(g.h!iwj)trem(ih..) (;) im(h)mén(h)sae(h') ma(g)je(hi)stá(ivHG)tis :(g.) (::) 
Ve(gi)ne(h)rán(j')dum(i) tu(h')um(h) ve(hi)rum,(h.) (;) et(h) ú(h')ni(g)cum(hi) Fí(i)li(hg)um :(g.) (::) 
San(g.h!iwj)ctum(i_[oh:h]h) quo(hi)que(h.) (;) Pa(h)rá(h')cli(g)tum(hi) Spí(i)ri(hg)tum.(g.) (::) 
Tu(gh) Rex(h) gló(hg)ri(h)ae,(hi) (,) Chri(g.)ste.(e.) (::) 
Tu(e) Pa(gh)tris(h.) (,) sem(h)pi(h)tér(hg)nus(h) es(hi) Fí(g)li(g)us.(e.) (::) 
<alt>Fit reverentia</alt>
Tu(e) ad(gh~) li(h)be(h)rán(h)dum(h) sus(h)ce(h)ptú(h')rus(g) hó(h!iwj)mi(ivHG)nem,(h.) (;) non(e) hor(gh)ru(h)í(h)sti(h) Vír(hg)gi(h)nis(hi) ú(g)te(g)rum.(e.) (::) 
Tu(e) de(gh)ví(h)cto(h) mor(h)tis(h') a(g)cú(h!iwj)le(ivHG)o,(h.) (;) a(e)pe(gh)ru(h)í(h)sti(h) cre(h)dén(i)ti(h)bus(h) re(hg)gna(h) cae(hi)ló(g.)rum.(e.) (::) 
Tu(e) ad(gh~) déx(h)te(h)ram(h) De(h')i(g) se(h!iwj)des,(ih..) (;) in(h) gló(hg)ri(h)a(hi) Pa(g.)tris.(e.) (::) 
Ju(gh)dex(h) cré(hi)de(h)ris(h.) (;) es(hg)se(h) ven(hi)tú(g.)rus.(e.) (::) 
<alt>Sequens versus dicitur flexis genibus</alt> 
<nlba>Te(e) er(gh~)go(h) quaé(h)su(h)mus,(h.) (,) tu(h)is(h) fá(h)mu(h')lis</nlba>(g) súb(h!iwj)ve(ivHG)ni,(h.) (;) quos(e) pre(gh)ti(h)ó(h)so(h) sán(i)gui(h)ne(hg) red(h)e(hi)mí(g.)sti.(e.) (::) 
Æ(f)tér(fe)na(df) fac(evDC.) (;) cum(c) san(e)ctis(f) tu(g)is(g) in(g) gló(g)ri(g)a(g) nu(h)me(f)rá(gvFE)ri.(e.) (::) 
Sal(c)vum(df) fac(f) pó(f)pu(f)lum(f) tu(f)um(f) Dó(d)mi(f)ne,(evDC.) (;) et(c) bé(df)ne(f)dic(f) he(f)re(fe)di(d)tá(f)ti(gh) tu(gvFE)ae.(e.) (::) 
Et(f) re(fe)ge(d) e(f)os,(evDC.) (;) et(c) ex(df)tól(fe)le(d) il(g)los(g.) (,) us(g)que(g) in(h) ae(f)tér(gvFE)num.(e.) (::) 
Per(e) sín(gh)gu(h)los(hg) di(h!iwj)es,(ih..) (;) be(h)ne(hg)dí(h)ci(hi)mus(g.) te.(e.) (::) 
<alt>Fit reverentia, secundum consuetudinem</alt>
Et(e) lau(gh~)dá(h)mus(h) no(h)men(h) tu(h)um(h') in(g) saé(h!iwj)cu(ivHG)lum,(h.) (;) et(h) in(h) saé(hg)cu(h)lum(hi) saé(g)cu(g)li.(e.) (::) 
Di(e)gná(gh)re(h) Dó(h)mi(h)ne(h) di(h')e(g) i(h!iwj)sto(ih..) (;) si(e)ne(gh) pec(h)cá(h)to(h) nos(hg) cu(h)sto(hi)dí(g.)re.(e.) (::) 
Mi(e)se(gh)ré(h)re(h) no(h')stri(g) Dó(h!iwj)mi(ivHG)ne,(h.) (;) mi(h)se(hg)ré(h)re(hi) no(g.)stri.(e.) (::) 
Fi(e)at(gh) mi(h)se(h)ri(h)cór(h)di(h)a(h) tu(h)a(h.) Dó(h)mi(h')ne(g) su(h!iwj)per(ivHG) nos,(h.) (;) quem(e)ád(gh~)mo(h)dum(h) spe(h)rá(hg)vi(h)mus(hi) in(g.) te.(e.) (::) 
In(c) te(df) Dó(fe)mi(d)ne(f) spe(gh)rá(gf/gf)vi :(e[ll:1]d..) (;) non(c) con(df)fún(fe)dar(dc) in(f) ae(gh)tér(gfgvF~E~)num.(e.) (::)}
\vspace{7mm}
\normalsize

% Whereas \sloppy will loosen the alignment constraints on the whole document, the sloppypar environment can limit the effect of the damage to a localized area.
%\sloppy
%
%\begin{sloppypar}
%\begin{spacing}{0.9}
\begin{paracol}{3}
\switchcolumn[0]
		Ciebie Boże chwalimy:\\Ciebie Panem wyznawamy.\\Tobie wiecznemu Ojcu,\\cała ziemia cześć oddawa.\\Tobie wszyscy Anieli,\\Tobie niebiosa i wszystkie mocarstwa,\\Tobie Cherubim i Serafim\\nieustannym głosem śpiewają:\\Święty, Święty, Święty,\\Pan Bóg zastępów.\\Pełne są niebiosa i ziemia\\Majestatu chwały Twojéj.\\Ciebie przesławny chór Apostolski,\\Ciebie Proroków chwalebny poczet,\\Ciebie świetne wojsko \\Męczenników chwali.\\Ciebie po całym okręgu ziemi,\\święty Kościół wyznawa,\\Ojca niezmiernego Majestatu,\\Chwalebnego téż i prawdziwego\\
	\switchcolumn[1]
	i jedynego Syna Twego,\\I Ducha świętego Pocieszyciela.\\Tyś królem chwały, Chryste.\\Tyś Ojca przedwiecznym jest Synem.\\Ty na odkupienie, \\chcąc przyjąć człowieczeństwo,\\nie pogardziłeś panieńskim żywotem.\\Ty skruszywszy żądło śmierci, \\otworzyłeś wierzącym królestwo niebieskie.\\Ty na prawicy Bożéj siedzisz\\we chwale Ojca.\\Sądzić nas kiedyś, \\wierzymy, że przyjdziesz.\\Przetoż Cię prosimy, \\racz wspomódz sługi Twoje,\\które przenajdroższą krwią odkupiłeś.\\Spraw, abyśmy byli z Świętymi Twy-mi\\
	\switchcolumn[2]
	do wiecznéj chwały policzeni.\\Racz zbawić lud Twój Panie \\i błogosław dziedzictwu Swojemu.\\I zarządzaj niém, \\podwyższając je aż na wieki.\\Każdego dnia błogosławimy Ciebie.\\I chwalimy Imię Twoje \\na czas wieczny, i na wieki wieków.\\Racz, Panie, dnia tego (téj nocy) \\od grzechu nas zachować.\\Zmiłuj się nad nami, Panie, \\Zmiłuj się nad nami.\\Niech będzie miłosierdzie Twoje,\\Panie nad nami,\\gdy pokładamy nadzieję w Tobie.\\W Tobiem Panie, nadzieję położył, \\niech nie będę zawstydzon na wieki.\footnote{tłum. z "Ołtarzyk polski katolickiego nabożeństwa", wydanie czwarte pomnożone, Grudziądz 1883}
\end{paracol}	
%\end{sloppypar}
%\end{spacing}
\end{document}